\section{Introducción y Objetivo del Sistema}
\label{sec:introduccion_objetivos}

\subsection{Contextualización de la Problemática}
En el ecosistema económico de El Salvador, las micro, pequeñas y medianas empresas (MIPYMES) del sector gastronómico, cafeterías de especialidad y panaderías artesanales constituyen un pilar fundamental del empleo y la actividad comercial urbana. No obstante, la mayoría de estos establecimientos operan con una infraestructura administrativa altamente fragmentada.

La práctica operativa habitual se basa en el uso simultáneo de cuadernos de notas para la toma de comandas, hojas de cálculo aisladas para registrar compras y planillas, y terminales de cobro independientes o sistemas de punto de venta desconectados de la contabilidad de costos. Esta desconexión genera fallas críticas:
\begin{itemize}
    \item \textbf{Desfase Crítico en Inventarios:} La venta de productos elaborados (por ejemplo, una taza de café cappuccino o un croissant de almendras) no actualiza automáticamente las existencias de leche, granos de café o mantequilla, provocando desabastecimientos imprevistos o sobrecompras de insumos perecederos.
    \item \textbf{Vulnerabilidad en Arqueos de Caja:} La ausencia de un protocolo formal de arqueo ciego y turnos cerrados dificulta identificar faltantes o sobrantes de efectivo al final de la jornada.
    \item \textbf{Inexactitud en el Costeo de Productos:} Sin un control riguroso del Costo Promedio Ponderado de cada insumo entrante, los precios de venta al público se fijan por intuición y no con base en márgenes de contribución reales.
    \item \textbf{Complejidad Tributaria y Laboral:} El cumplimiento de la normativa fiscal salvadoreña (desglose del 13\% de IVA en Factura de Consumidor Final y Comprobante de Crédito Fiscal) y los cálculos laborales (retenciones de ISSS, AFP y tramos de Impuesto sobre la Renta) suelen delegarse en procesos manuales propensos a error humano y sanciones administrativas.
\end{itemize}

\subsection{Objetivo General}
Diseñar, desarrollar e implementar \textbf{artesa}, un sistema unificado de Planificación de Recursos Empresariales (ERP) y Punto de Venta (POS) con arquitectura modular web y persistencia relacional en PostgreSQL 16, orientado a optimizar la gestión operativa, el control de inventario basado en recetas y la conformidad fiscal y laboral de cafeterías y panaderías artesanales en El Salvador.

\subsection{Objetivos Específicos}
\begin{enumerate}
    \item \textbf{Unificación Transaccional del Inventario:} Automatizar la deducción de existencias en tiempo real mediante funciones PL/pgSQL atómicas que descompongan los productos vendidos en sus insumos elementales según recetas preconfiguradas.
    \item \textbf{Operación Ágil de Mostrador y Cocina:} Proveer una interfaz de punto de venta reactiva y una pantalla de visualización de comandas en cocina (Kitchen Display System - KDS) para sincronizar pedidos, preparación y despacho.
    \item \textbf{Blindaje Financiero de Caja:} Implementar una máquina de estados finitos para la gestión de turnos de caja con arqueo ciego, registro de ingresos/egresos extraordinarios y auditoría de discrepancias.
    \item \textbf{Control de Abastecimiento y Cuentas por Pagar:} Gestionar el catálogo de proveedores locales, registro de facturas de compra con actualización de costo promedio y seguimiento de obligaciones pendientes (CxP).
    \item \textbf{Cálculo de Planilla Paramétrico:} Automatizar la liquidación de sueldos conforme a los tramos de renta de El Salvador y las tasas legales de cotización de seguridad social (ISSS, AFP, INSAFORP).
    \item \textbf{Auditoría y Trazabilidad Forense:} Registrar en bitácora inmutable las operaciones críticas del sistema, almacenando el estado anterior y posterior en formato JSON junto con la dirección IP y el agente de usuario del operador.
\end{enumerate}
